\documentclass[11pt]{article}
\usepackage{mathblatt}
% Prüfungsblatt nach mathe-nachhilfe pruefung.md § 5 (Runde Dreiecke, 11.10.2026), Hand-Satz; Präambel des Setzers.
% Präambel des Setzers (werkzeuge/setzer.py), wortgleich aus dem Hand-Satz T6B (bau/proben/2026-10-09/kathete-berechnen), dort aus K4W.
\makeatletter
\setlength{\mbpfbe}{0mm}% keine Punkte (Bauregel 6.4)
% eingerückt (Bauregel 4.7, 6.6): \gEin Einzug, \gSchrift eine Stufe kleiner
\newlength{\gEin}\setlength{\gEin}{0pt}
\let\gSchrift\relax
\renewcommand{\pfkreuzzeile}[1]{\par\noindent\hangindent3.2em\hangafter1\hspace*{1.6em}$\square$\ #1\par}
\newcommand{\gkreuz}[1]{$\square$\ #1\hspace{2em}}
% Bauregel 6.7: links, was man ansieht, rechts, was man tut; Spaltengrenze überall an derselben Stelle
\newsavebox{\gbox}\newlength{\gLinks}
\newcommand{\gzwei}[2]{\par\smallskip\noindent\sbox{\gbox}{#1}%
  \setlength{\gLinks}{\dimexpr0.5\textwidth-\gEin-\mbpfnr\relax}%
  \ifdim\wd\gbox>\dimexpr\gLinks-4mm\relax
    \usebox{\gbox}\par\smallskip\noindent\begin{minipage}{\linewidth}#2\end{minipage}%
  \else
    \begin{minipage}[t]{\gLinks}\vspace{0pt}\usebox{\gbox}\end{minipage}%
    \begin{minipage}[t]{\dimexpr\linewidth-\gLinks\relax}\vspace{0pt}\raggedright #2\end{minipage}%
  \fi\par}
\RenewDocumentEnvironment{pfaufg}{m m m}{%
  \begin{lrbox}{\mbaufgabebox}\begin{minipage}[t]{\linewidth}%
  \gSchrift\noindent\hspace*{\gEin}\mbpfrandmarke{#1}%
  \begin{minipage}[t]{\mbpfnr}\strut\textbf{#2}\end{minipage}%
  \begin{minipage}[t]{\dimexpr\linewidth-\gEin-\mbpfnr\relax}\raggedright\strut\ignorespaces}%
 {\end{minipage}%
  \end{minipage}\end{lrbox}%
  \setlength{\mbaufgabehoehe}{\dimexpr\ht\mbaufgabebox+\dp\mbaufgabebox\relax}%
  \par\addvspace{7pt}\Needspace*{\mbaufgabehoehe}%
  \noindent\usebox{\mbaufgabebox}\par\addvspace{7pt}}
\makeatother
% Rechenraum als Karo (Bauregel 6.9): n Rechenschritte = n mal 10 mm
\newcommand{\karo}[1]{\par\smallskip\noindent\begin{tikzpicture}
  \clip (0,0) rectangle ({\linewidth-1mm},{#1*10mm});
  \draw[black!16,line width=0.3pt,step=5mm] (0,0) grid ({\linewidth},{#1*10mm});
  \end{tikzpicture}\par}
\newcommand{\kk}[1]{$\square$\,#1\hspace{0.9em}}
\newcommand{\linie}[1]{\,{\color{black!45}\rule[-1pt]{#1}{0.4pt}}\,}
% Zeile am Ende jedes Durchgangs: Originale klein grau, darüber; dann Vorgänger/Nachfolger (Bauregel 3.4, 6.5)
\newcommand{\blattende}[2]{\par\vfill\noindent\begin{minipage}{\linewidth}{\scriptsize\color{mbgrau}Originale: #1\par}\smallskip
{\footnotesize #2\par}\end{minipage}}
\newcommand{\pkt}{\textperiodcentered{}}

\newcommand{\kennung}{LXA}

\begin{document}
\pfheftstil
\pfheftkopf[\kennung]{Seite im allgemeinen Dreieck, Sinussatz}{P10}

% ===================== Stufe 1 =====================
\Needspace*{0.25\textheight}\pfstufekopf{Dritten Winkel finden}{geprüft 2025 \pkt{} 2024}

% Bank trigonometrie-e4-k1-s5-v7 (Nebenwinkel, dann Winkelsumme; nichts mit Sinus)
\begin{pfaufg}{}{1.}{}
Die Punkte $A$, $B$ und $D$ liegen auf einer Geraden. Berechne die Winkel $\beta$ und $\gamma$ im Dreieck $ABC$.
\gzwei{\begin{tikzpicture}[font=\small]\draw[line width=0.8pt] (0.000,0.000) -- (1.821,0.000) -- (2.622,2.200) -- cycle;\draw[line width=0.8pt] (1.821,0.000) -- (3.441,0.000);\draw[line width=0.5pt] (0.420,0.000) arc (0.00:40.00:0.42);\node at (0.651,0.237) {$40^\circ$};\draw[line width=0.5pt] (2.241,0.000) arc (0.00:70.00:0.42);\node at (2.386,0.396) {$70^\circ$};\node[anchor=north,inner sep=2pt] at (0.911,-0.050) {5\,cm};\node[anchor=north east,inner sep=2pt] at (0.000,0.000) {$A$};\node[anchor=north,inner sep=2pt] at (1.821,0.000) {$B$};\node[anchor=south,inner sep=2pt] at (2.622,2.200) {$C$};\node[anchor=north,inner sep=2pt] at (3.441,0.000) {$D$};\end{tikzpicture}}{%
\karo{2}}
\fusshilfe{\mbox{1:~$\beta = 110^\circ$, $\gamma = 30^\circ$}}
\end{pfaufg}

% 2024-OS-K6d, herausgelöst (nur der Winkel bei C; Länge BC steht in der Fundstellentabelle), eigener Wortlaut
\begin{pfaufg}{nach~P10\\’24~\pkt{}~6d}{2.}{}
Eine Seilbahn führt vom Ort $B$ im Tal zum Ort $A$ auf einem Berg. Eine zweite Seilbahn verbindet $B$ mit dem Ort $C$ auf einem anderen Berg. Bei $A$ ist $\alpha = 38^\circ$, bei $B$ ist $\beta_2 = 108^\circ$. Berechne den Winkel bei $C$ im Dreieck $ABC$.
\gzwei{\begin{tikzpicture}[line width=0.7pt,scale=0.62]
\coordinate (A) at (0,2.87);\coordinate (F) at (0,0);\coordinate (B) at (2.55,0);\coordinate (C) at (6.42,1.69);
\draw (A) -- (F) -- (B) -- cycle;\draw (B) -- (C) -- (A);\mbrwmarke{(F)}{(B)}{(A)}
\draw[line width=0.5pt] ($(A)+(-48.4:0.55)$) arc (-48.4:-10.4:0.55);\node[font=\small] at ($(A)+(-29:0.95)$) {$\alpha$};
\draw[line width=0.5pt] ($(B)+(131.6:0.5)$) arc (131.6:180:0.5);\node[font=\small] at ($(B)+(156:0.95)$) {$\beta_1$};
\draw[line width=0.5pt] ($(B)+(23.6:0.4)$) arc (23.6:131.6:0.4);\node[font=\small] at ($(B)+(78:0.75)$) {$\beta_2$};
\node[font=\small,anchor=south] at (A) {$A$};\node[font=\small,anchor=north east] at (F) {$F$};\node[font=\small,anchor=north] at (B) {$B$};\node[font=\small,anchor=south west] at (C) {$C$};
\end{tikzpicture}}{%
\karo{1}}
\fusshilfe{\mbox{2:~$34^\circ$}}
\end{pfaufg}

% ===================== Stufe 2 =====================
\Needspace*{0.25\textheight}\pfstufekopf{Seite berechnen}{geprüft 2025 \pkt{} 2024}

% Bank trigonometrie-e4-k1-s0-v6 (vollständiges Paar finden, nichts rechnen)
\begin{pfaufg}{}{3.}{}
Im Dreieck $ABC$ kennst du $b = 7$~cm, $\beta = 62^\circ$ und $\gamma = 49^\circ$. Kreuze an, welches Paar aus Seite und Gegenwinkel ganz bekannt ist.
\gzwei{\begin{tikzpicture}[font=\small]\draw[line width=0.8pt] (0.000,0.000) -- (3.362,0.000) -- (1.276,2.400) -- cycle;\draw[line width=0.5pt] (0.420,0.000) arc (0.00:62.00:0.42);\node at (0.591,0.355) {$\alpha$};\draw[line width=0.5pt] (3.087,0.317) arc (131.00:180.00:0.42);\node at (2.696,0.304) {$\beta$};\draw[line width=0.5pt] (1.079,2.029) arc (-118.00:-49.00:0.42);\node at (1.366,1.608) {$\gamma$};\node[anchor=north,inner sep=2pt] at (1.681,-0.050) {$c$};\node[anchor=south west,inner sep=2pt] at (2.357,1.233) {$a$};\node[anchor=south east,inner sep=2pt] at (0.594,1.223) {$b$};\node[anchor=north east,inner sep=2pt] at (0.000,0.000) {$A$};\node[anchor=north west,inner sep=2pt] at (3.362,0.000) {$B$};\node[anchor=south,inner sep=2pt] at (1.276,2.400) {$C$};\end{tikzpicture}}{%
\pfkreuzab{A}{$b$ und $\gamma$}\smallskip\pfkreuzab{B}{$b$ und $\beta$}\smallskip\pfkreuzab{C}{$c$ und $\gamma$}}
\fusshilfe{\mbox{3:~B}}
\end{pfaufg}

% Bank trigonometrie-e4-k1-s1-v6, e4-k1-s3-v4 (Päckchen: Paar bekannt; Paar erst nach der Winkelsumme)
\begin{pfaufg}{}{4.}{}
Berechne die gesuchte Seite. Runde auf eine Stelle nach dem Komma.\par\smallskip
\gzwei{\parbox[t]{\dimexpr0.5\textwidth-\gEin-\mbpfnr-5mm\relax}{\raggedright \textbf{a)}\ $a = 8$~cm, $\alpha = 70^\circ$, $\beta = 45^\circ$; gesucht: $b$}}{\karo{2}}
\gzwei{\parbox[t]{\dimexpr0.5\textwidth-\gEin-\mbpfnr-5mm\relax}{\raggedright \textbf{b)}\ $c = 7$~cm, $\alpha = 48^\circ$, $\beta = 64^\circ$; gesucht: $a$}}{\karo{2}}
\fusshilfe{\mbox{4:~a) 6,0 cm; b) 5,6 cm}}
\end{pfaufg}

% Bank trigonometrie-e4-k1-s4-v4 (Sache, Nebenwinkel, stumpfer Winkel im Paar)
\begin{pfaufg}{}{5.}{}
Eine Zeltstange $\overline{BC}$ steht schräg. Ein Seil $s$ hält sie vom Hering $A$ aus. Die Punkte $A$, $B$ und $D$ liegen auf einer Geraden. Wie lang ist das Seil?
\gzwei{\begin{tikzpicture}[font=\small]\draw[line width=0.8pt] (0.000,0.000) -- (1.802,0.000) -- (2.519,1.968) -- cycle;\draw[line width=0.8pt] (1.802,0.000) -- (3.600,0.000);\draw[line width=0.5pt] (0.420,0.000) arc (0.00:38.00:0.42);\node at (0.683,0.235) {$38^\circ$};\draw[line width=0.5pt] (2.222,0.000) arc (0.00:70.00:0.42);\node at (2.368,0.396) {$70^\circ$};\node[anchor=north,inner sep=2pt] at (0.901,-0.050) {2{,}5\,m};\node[anchor=south east,inner sep=2pt] at (1.229,1.023) {$s$};\node[anchor=north east,inner sep=2pt] at (0.000,0.000) {$A$};\node[anchor=north,inner sep=2pt] at (1.802,0.000) {$B$};\node[anchor=south,inner sep=2pt] at (2.519,1.968) {$C$};\node[anchor=north,inner sep=2pt] at (3.600,0.000) {$D$};\end{tikzpicture}}{%
\karo{3}}
\fusshilfe{\mbox{5:~4,4 m}}
\end{pfaufg}

% 2025-OS-K4c, eigener Wortlaut (Sternchenaufgabe; Rampe vorn aus 4a, 4b weggelassen)
\begin{pfaufg}{nach~P10\\’25~\pkt{}~4c\,$\star$}{6.}{}
Am Hinterausgang soll eine mobile Rampe $y$ zwei Stufen überwinden. Vom Fuß der Rampe bis zum Fuß der Treppe sind es waagerecht $160$~cm. Die Rampe steigt unter $5^\circ$ an. Die gestrichelte Linie vom Fuß der Treppe zur oberen Stufenkante bildet mit dem Boden einen Winkel von $141^\circ$. Berechne die Länge $y$ der Rampe.
\gzwei{\begin{tikzpicture}[line width=0.7pt,scale=1.15]
\coordinate (R) at (0,0);\coordinate (T) at (4,0);\coordinate (S) at (4.93,0.75);
\fill[black!15] (4,0) -- (4.465,0) -- (4.465,0.375) -- (4.93,0.375) -- (4.93,0.75) -- (5.5,0.75) -- (5.5,0) -- cycle;
\draw (R) -- (5.5,0);\draw[line width=1.1pt] (R) -- (S);\draw[dashed] (T) -- (S);
\draw[line width=0.5pt] (1.6,0) arc (0:8.65:1.6);\node[font=\scriptsize,anchor=west,inner sep=1pt] at (1.66,0.11) {$5^\circ$};
\draw[line width=0.5pt] ($(T)+(38.9:0.3)$) arc (38.9:180:0.3);\node[font=\small,anchor=east] at ($(T)+(170:0.32)$) {$141^\circ$};
\node[font=\small,anchor=north] at (2.6,-0.05) {$160$ cm};\node[font=\small,anchor=south east] at (2.4,0.37) {$y$};
\end{tikzpicture}}{%
\karo{3}}
\fusshilfe{\mbox{6:~180,1 cm}}
\end{pfaufg}

\pfblattende{11}{26 & 25 & 24 & 23 & 22 & 21 & 20 & 19 & 18 & 17 & 16}%
{ & 4c & 6d & & & 3c & 7c & 3c & 4c & 4c & }%
{Weiter: Symmetrie}
\end{document}
